\documentclass{article}
\usepackage{amsmath}
\usepackage{amsfonts}
\usepackage{amssymb}
\usepackage{amsthm}
\usepackage[margin=2cm]{geometry}

% No paragraph indentation
\setlength{\parindent}{0pt}

% 2 lines of spacing between paragraphs
\setlength{\parskip}{1\baselineskip}

% Definition-style (plain vs definition style)
\theoremstyle{definition}
\newtheorem{definition}{Definition}[section]

\title{A short note on cubic spline interpolation}
\author{Quasar Chunawala}
\begin{document}
\maketitle

\section{Introduction}

We start by formally defining a spline function of degree $k - 1 \geq 0$. 

\begin{definition}
    A spline function $S(x)$ of degree $k - 1 \geq 0$ on a grid: 
    $$
    \Delta = \{a = x_0 < x_1 < \ldots < x_n = b\}
    $$
    of distinct knots is a real function with the following properties:
    \begin{enumerate}
        \item For $x \ in [x_i, x_{i+1}]$, $S(x)$  is a polynomial of degree $< k$.
        \item Its first $k - 2$ derivatives are continuous, that is $S(x) \in C^{k-2}[a,b]$.
    \end{enumerate}
\end{definition}

The space of all spline functions of degree $k - 1 \geq 0$ is denoted by $S_{\Delta,k}$. From the definition, it follows that if $S_1(x)$ and $S_2(x)$ are spline functions, so is $c_1 S_1(x) + c_2 S_2(x)$. Clearly, $S_{\Delta, k}$ is a linear vector space. 

\section{Cubic Spline Interpolant}

A cubic spline uses cubic polynomials of degree $3$ to interpolate between successive knots. A subic spline interpolant $S(x)$ satisfies the following conditions:

1. Each piece or segment $S_j(x)$ on the interval $[x_j, x_{j+1}]$ is a cubic polynomial.

2. $S_j(x_j) = y_j$ for each $j = 0, 1, \ldots, n$.

3. $S_j(x_{j+1}) = S_{j+1}(x_{j + 1})$, $S_j'(x_{j+1}) = S_{j+1}'(x_{j+1})$ and $S_j''(x_{j+1}) = S_{j+1}''(x_{j+1})$.

Condition (3) is just a translation of the requirement that the whole spline should be $k-2$ times differentiable and these derivatives must be continuous. As a crude approximation, we know that a continuous curve is one that can be drawn without lifting pen from paper. The individul segments of $S(x)$ are cubics, so they are continuous themselves. The only consideration is the continuity at knot points.

For a curve to be continous at any given knot-point $c$, whether we approach it from the left or we approach it from the right, the spline funtion $S(x)$ must approach $S(c)$ regardless. Suppose the left piece is $S_j$ and the right piece is $S_{j+1}$. This motivates $S_j(x_{j + 1}) = S_{j + 1}(x_{j + 1})$. A similar argument follows for the first and second derivatives because we would like those to be continous.

\section{Construction of the Cubic Spline}


 
\end{document}